\begin{theorem}[Vertical movement with a common tree background]
    \label{thm:peps_torus_gauged_vertical_flux_move}
    \lean{TNLean.PEPS.IsGIsometric.exists_unitary_torusGaugedVerticalFluxMove,
        TNLean.PEPS.torusGaugedVerticalFluxAssignment,
        TNLean.PEPS.torusGaugedVerticalFluxAssignment_one}
    \leanok
    \uses{thm:peps_torus_vertical_flux_move,
        thm:peps_regular_gauged_cycle_permutation,
        thm:peps_regular_two_cycle_flux_move}
    Let the original torus site tensor be $G$-isometric for the regular action
    of a finite group $G$. Assume width at least three and height at least
    four. At any position $v$, use the actual two-by-three six-site region
    $R$ and its derived path tree. Write $u_0(g),u_1(g)$ for the lower-chord
    and double-chord insertions of the preceding vertical construction.
    They retain the ordered seam element $g^{-1}$ when the rightward step
    crosses the horizontal seam. For a vertex gauge $k:R\to G$, reconstruct
    the internal assignment as
    \begin{align}
        s_{k,b,g}(e)
        =k_{\mathrm{head}(e)}u_b(g)(e)k_{\mathrm{tail}(e)}^{-1},
        \qquad b\in\{0,1\}.
    \end{align}
    There is one six-spin unitary $W$, chosen before every $k$, direction,
    $g$, boundary configuration $\theta$, and literal exterior and crossing
    assignment $u$, such that
    \begin{align}
        W\Psi_R(s_{k,0,g},u,\theta)     & =\Psi_R(s_{k,1,g},u,\theta), \\
        W^*\Psi_R(s_{k,1,g},u,\theta)   & =\Psi_R(s_{k,0,g},u,\theta), \\
        (W\otimes I)\Psi(s_{k,0,g},u)   & =\Psi(s_{k,1,g},u),          \\
        (W^*\otimes I)\Psi(s_{k,1,g},u) & =\Psi(s_{k,0,g},u).
    \end{align}
    Both sides retain the same $k$ and the same exterior and crossing $u$.
    The identity gauge recovers the literal vertical insertions. The
    assertions include both periodic seams and concern actual original-spin
    contractions. This is a finite-torus common-background specialization of
    \cite[accessible systems and Theorem~6.16]{Schuch2010PEPS}; it does not
    compare distinct crossing presentations or establish a prescribed
    four-endpoint braid. See \cite{gap:rmp_peps_quantum_double_g_isometry}.
\end{theorem}

\begin{proof}\leanok
    Recover the literal insertions as the derived cycle assignments and
    reconstruct them using the common vertex gauge. The same gauge and
    crossing operators give the same actual boundary transport. Apply the
    uniform physical cycle permutation to these normalized columns. The
    actual cut expansion gives the global identity, and unitarity gives
    both reverse identities by multiplication with the adjoint.
\end{proof}
